\BookPart{From molecules to genomes}{genomes}{%
Living cells compute with DNA every second: genes are switched on and off by
regulatory proteins, cells remember their history, and development unfolds a
genome into a body. Part~VIII connects engineered molecular programs to this
biological computation. It models gene regulation as a dynamical system, asks
what it takes to run a molecular program in serum or inside a cell, and closes
by asking precisely which ideas carry from molecular computation to genomic
computation, and where the analogy fails.

\textbf{Chapters in this part}
\begin{itemize}\setlength{\itemsep}{2pt}
\item \textbf{30\quad Genomic organisation, regulation and development.} How do regulatory networks create stable cell states and memories?
\item \textbf{31\quad Modern molecular programming and evidence.} How do DNA programs behave in realistic environments, and how is that behaviour demonstrated?
\item \textbf{32\quad From molecular computation to genomic computation.} Which properties of a molecular state does a representation preserve, and what does that mean for genomic models?
\end{itemize}}
